\documentclass[11pt,a4paper]{amsart}
\usepackage[T1]{fontenc}
\usepackage{lmodern}
\usepackage{amssymb}
\usepackage{microtype}
\usepackage[hyphens]{url}
\usepackage[colorlinks=true,linkcolor=black,citecolor=black,urlcolor=black]{hyperref}
\usepackage[margin=2.5cm]{geometry}

\hypersetup{
  pdftitle={Two-colourable planes over number fields},
  pdfauthor={Sergi Gal\'an},
}

\newtheorem{theorem}{Theorem}
\newtheorem{proposition}[theorem]{Proposition}
\newtheorem{lemma}[theorem]{Lemma}
\newtheorem{corollary}[theorem]{Corollary}
\theoremstyle{definition}
\newtheorem{question}[theorem]{Question}
\newtheorem{remark}[theorem]{Remark}
\newtheorem{example}[theorem]{Example}

\newcommand{\F}{\mathbb{F}}
\newcommand{\Q}{\mathbb{Q}}
\newcommand{\R}{\mathbb{R}}
\newcommand{\Z}{\mathbb{Z}}
\newcommand{\Oo}{\mathcal{O}}
\newcommand{\m}{\mathfrak{m}}
\newcommand{\N}{\operatorname{N}}
\newcommand{\chiB}{\chi_{\mathrm{B}}}
\newcommand{\chiM}{\chi_{\mathrm{M}}}
\newcommand{\chiloc}{\chi_{\mathrm{loc}}}

\title{Two-colourable planes over number fields}
\author{Sergi Gal\'an}
\thanks{These results were obtained with the assistance of an AI system (Claude, by Anthropic), which carried
out the computations and the checks and helped to draft this note. Two further sessions of the same system,
given only the manuscript, refereed it; no mathematician has checked it yet. The author is not a professional
mathematician and submits it for independent verification.}
\date{Draft, 3 October 2026}
\subjclass[2020]{Primary 05C15; Secondary 11R04, 11S15, 52C10}
\keywords{Hadwiger--Nelson problem, unit-distance graph, number field, ramification, local--global principle}

\begin{document}

\begin{abstract}
For a number field $F$, let $\chi(F^2)$ be the chromatic number of the graph on $F^2$ in which $(x,y)$ and
$(x',y')$ are adjacent when $(x-x')^2+(y-y')^2=1$. We show that $\chi(F^2)=2$ if and only if some prime of $F$
above $2$ ramifies in $F(i)$. That such a prime gives a $2$-colouring, by reduction, was known; we prove the
converse: if no prime above $2$ ramifies in $F(i)$, the graph has an odd cycle. The proof combines Fischer's
observation that a $2$-colouring is a ring homomorphism with Chevalley's extension theorem for valuation rings.
The criterion contains the known cases of real quadratic and real multiquadratic fields and of fields of odd
degree, and it holds for the plane of any nondegenerate binary form $x^2+cy^2$, with $F(\sqrt{-c})$ in place of
$F(i)$. We then ask whether, for a real number field, $\chi(F^2)$ is always the least number of colours of a
colouring of a single completion $F_v^2$ that is locally constant (Borel at a real place); the theorem is the case
of two colours. For three colours the question
predicts $\chi(\Q(\sqrt d)^2)\ge4$ for every squarefree $d\equiv11\pmod{12}$, and for four colours a positive
answer for $\Q(\sqrt{167})$ would give a triangle-free $5$-chromatic unit-distance graph in the plane.
\end{abstract}

\maketitle

\section{Introduction}

The Hadwiger--Nelson problem asks for the chromatic number $\chi(\R^2)$ of the unit-distance graph of the plane;
it is $5$, $6$ or $7$ \cite{deGrey,EI,Soifer}. For a subfield $F\subset\R$ the plane $F^2$ is a subgraph, and
since the adjacency $(x-x')^2+(y-y')^2=1$ is algebraic, $\chi(F^2)$ makes sense for every field $F$. Woodall
showed $\chi(\Q^2)=2$ \cite{Woodall}. For real quadratic fields, $\chi(\Q(\sqrt N)^2)=2$ exactly when
$N\not\equiv3\pmod 4$: Johnson proved that these planes are $2$-colourable \cite{Johnson1987} (we know his paper
through \cite{Fischer1990,Payne}), and Fischer proved both directions \cite[Theorem~8]{Fischer1990}. Moorhouse
gave proofs of the $2$-colourable cases and showed $\chi(F^2)=2$ for every field of odd degree
\cite[Theorem~7.1, Lemma~8.4, Theorem~8.5]{Moorhouse}.

The colourings of Woodall, Fischer and Moorhouse are reductions modulo a prime above $2$. Madore's reduction
argument \cite[Proposition~3.2, \P6.6]{Madore} needs $x^2+y^2$ to stay anisotropic modulo the prime, which never
happens above $2$, where $x^2+y^2\equiv(x+y)^2$. For primes above $2$ he gave a version modulo the square of the
prime, which applies when a prime above $2$ is unramified over $\Q$ \cite[\P3.7, Proposition~3.8]{Madore}, and a separate argument for
$\Q(\sqrt2)$ \cite[Proposition~3.9]{Madore}. Writing points as $z=x+iy\in F(i)$, a public repository of July
2026 \cite{MM} proved the general statement: if a prime of $F$ above $2$ ramifies in $F(i)$, reduction modulo
the prime above it gives a $2$-colouring \cite[Theorem~A$'$]{MM}. We recall its short proof as
Proposition~\ref{prop:if}, in the slightly stronger form of a locally constant colouring of the completion; the
argument is the same.

The other direction, an odd cycle whenever no prime above $2$ ramifies in $F(i)$, had been proved field by field.
For real quadratic fields Fischer proved it with the ring homomorphism attached to a $2$-colouring
\cite[Theorems~1, 7 and~8]{Fischer1990}, and Moorhouse, for $\Q(\sqrt p)$ with $p$ prime, with a cycle from Pell's
equation \cite[Theorem~8.6]{Moorhouse}; it follows for every field with such a quadratic subfield. We prove it for every number field, by combining
Fischer's ring homomorphism with Chevalley's extension theorem.

\begin{theorem}\label{thm:main}
Let $F$ be a number field. Then $\chi(F^2)=2$ if and only if some prime of $F$ above $2$ ramifies in $F(i)$.
\end{theorem}

The \emph{only if} direction, Proposition~\ref{prop:onlyif}, is the new part. We have not found it in the
literature we could consult, but we have not seen Johnson's paper \cite{Johnson1987}, his survey
\cite{Johnson2006} or the paper of Benda and Perles \cite{BP}. Section~\ref{sec:proof} proves the theorem,
Section~\ref{sec:cor} draws the consequences, and Section~\ref{sec:lg} puts it in a local--global frame.

\section{Proof}\label{sec:proof}

Let $L=F[t]/(t^2+1)$; it is the field $F(i)$ if $i\notin F$, and $F\times F$ if $i\in F$. Identify $F^2$ with
$L$ by $(x,y)\mapsto x+ty$; then $x^2+y^2$ is the norm $\N(z)=z\bar z$, and the unit vectors are
\[
T=\{z\in L:\ z\bar z=1\}.
\]
$T$ is a group under multiplication, so its additive span $A=\Z[T]$ is a subring of $L$. The graph is the Cayley
graph of the additive group $L$ with connection set $T$, so it is a disjoint union of translates of $A$, each a copy
of the Cayley graph $\mathrm{Cay}(A,T)$. Every $u\in T$ is a unit of $A$, since $u^{-1}=\bar u\in T$.

\begin{lemma}[Fischer {\cite[Theorem~1(iii)]{Fischer1990}}, for $F\subset\R$]\label{lem:ring}
$\chi(F^2)=2$ if and only if there is a ring homomorphism $\varphi:A\to\F_2$.
\end{lemma}

\begin{proof}
Since $-1\in T$, every element of $A$ is a sum of unit vectors. If $\mathrm{Cay}(A,T)$ is bipartite, let
$\varphi(a)$ be the parity of the length of a walk from $0$ to $a$; it is well defined and additive. It is also
multiplicative: if $a$ is a sum of $m$ unit vectors and $b$ a sum of $n$, then $ab$ is a sum of $mn$ products of
unit vectors, which are unit vectors. Conversely, a ring homomorphism sends each $u\in T$, a unit of $A$, to
$1\in\F_2^\times$, so $\varphi(a+u)=\varphi(a)+1$ and $\varphi$ is a proper $2$-colouring of $\mathrm{Cay}(A,T)$,
which extends to the translates. (The graph has edges, so $\chi\ge2$.)
\end{proof}

\begin{proposition}[{\cite[Theorem~A$'$]{MM}}; cf.\ {\cite[Propositions~3.2, 3.8 and~3.9]{Madore}}]\label{prop:if}
If a prime $v$ of $F$ above $2$ ramifies in $L$, then $F_v^2$, and hence $F^2$, has a $2$-colouring that is
locally constant for the topology of $F_v$.
\end{proposition}

\begin{proof}
Let $w$ be the place of $L$ above $v$. It is the only one, so conjugation preserves $|\cdot|_w$; a unit vector
$u$ of $F_v^2=L_w$ therefore satisfies $|u|_w^2=|u\bar u|_w=1$, so $u\in\Oo_w^\times$. Since $L_w/F_v$ is
ramified, the residue fields of $w$ and $v$ agree and conjugation acts trivially on them; reducing $u\bar u=1$
gives $u^2\equiv1\pmod{\m_w}$, and in characteristic $2$ this is $u\equiv1\pmod{\m_w}$. Choose an additive map
$\lambda$ from the residue field to $\F_2$ with $\lambda(1)=1$. Unit steps preserve the cosets of $\Oo_w$ in $L_w$;
fix a representative $r$ of each coset and colour $z$ by $\lambda\bigl((z-r)\bmod\m_w\bigr)$. A unit step changes
the colour. The colouring is constant on the cosets of $\m_w$.
\end{proof}

\begin{proposition}\label{prop:onlyif}
If there is a ring homomorphism $\varphi:A\to\F_2$, then some prime of $F$ above $2$ ramifies in $L$.
\end{proposition}

\begin{proof}
Suppose first that $i\in F$, so $L=F\times F$ and $T=\{(s,1/s):s\in F^\times\}$. For $s\neq0,1$ the sum of
$(s,1/s)$ and $(1-s,1/(1-s))$ is $(1,1/(s(1-s)))$; the difference of the sums for $s=2$ and $s=3$ is
$(0,-1/3)\in A$. The set $\{x\in F:(0,x)\in A\}$ is an ideal of the second projection of $A$, which contains
$F^\times$ and is therefore $F$; this ideal is non-zero, so $A\supseteq 0\times F$, by symmetry $A=F\times F$,
and $F\times F$ has no homomorphism to $\F_2$ because $2$ is invertible. So $i\notin F$ and $L$ is a field.

By Zorn's lemma the pairs $(B,\psi)$ of a subring $B\supseteq A$ of $L$ and a homomorphism $\psi$ from $B$ to an
algebraic closure of $\F_2$ extending $\varphi$ have a maximal element $(V,\psi)$, and by Chevalley's theorem
\cite[Lemma~5.19 and Theorem~5.21]{AM} $V$ is a valuation ring of $L$ with maximal ideal $\ker\psi$, which meets
$A$ in $\ker\varphi$. As $V$ contains $\Z$ and is integrally closed, $V\supseteq\Oo_L$; as $2\in\ker\varphi$,
$V\neq L$. So $V$ is the valuation ring of a place $w$ of $L$ above $2$, and every $u\in T$ satisfies
$u\equiv1\pmod{\m_w}$. Let $v$ be the place of $F$ below $w$.

The map $p\mapsto(p+t)/(p-t)$ is a homeomorphism from $F_v$ minus the roots of $p^2+1$ onto $T(F_v)\setminus\{1\}$,
where $T(F_v)$ is the group of norm-one elements of $L\otimes_F F_v$, and it maps $F$ onto $T\setminus\{1\}$; so
$T$ is dense in $T(F_v)$.
\begin{itemize}
\item If $v$ splits in $L$, then $L\otimes_F F_v=F_v\times F_v$ and $z\mapsto z_w$ identifies $T(F_v)$ with
$F_v^\times$, which is unbounded; so some element of $T$ is not integral at $w$, a contradiction.
\item If $v$ is inert in $L$, with residue field $\F_q$, then conjugation acts on the residue field $\F_{q^2}$ of
$w$ as $x\mapsto x^q$, and the unit vectors $y/\bar y$ with $y\in\Oo_w^\times$ reduce to $y^{1-q}$, which runs
through all $q+1\ge3$ roots of $x^{q+1}=1$. Each residue class is open, so some element of $T$ is not
$\equiv1\pmod{\m_w}$, a contradiction.
\end{itemize}
Hence $v$ ramifies in $L$.
\end{proof}

\begin{proof}[Proof of Theorem~\ref{thm:main}]
Combine Lemma~\ref{lem:ring} with Propositions~\ref{prop:if} and~\ref{prop:onlyif}.
\end{proof}

\section{Consequences}\label{sec:cor}

For a multiquadratic field the condition is a computation in the group $\Q_2^\times/\Q_2^{\times2}$, of order
$8$, generated by $-1$, $5$ and $2$, and written multiplicatively. For a subgroup $H$ let $\Q_2(\sqrt H)$ be the
field generated over $\Q_2$ by square roots of representatives of the elements of $H$; then $\Q_2(\sqrt5)$ is the
only unramified quadratic extension of $\Q_2$.

\begin{corollary}\label{cor:multi}
Let $F=\Q(\sqrt{a_1},\dots,\sqrt{a_k})$ with nonzero integers $a_j$. Then $\chi(F^2)=2$ unless some quadratic
subfield $\Q(\sqrt c)$, with $c\neq1$ a squarefree integer, has $c\equiv3\pmod4$; that is, $F^2$ is
$2$-colourable exactly when the plane of every quadratic subfield is.
\end{corollary}

For real $F$ this is not new: the \emph{if} half is \cite[Corollary~B]{MM}, and the \emph{only if} half follows
from Fischer's theorem, since $F^2$ contains the plane of each subfield. For non-real $F$ the same reasoning works
with Theorem~\ref{thm:main} in place of those results. We include the corollary to show how Theorem~\ref{thm:main}
reads in this case.

\begin{proof}
The primes of $F$ above $2$ are conjugate. Let $G$ be the image of $\langle a_1,\dots,a_k\rangle$ in
$\Q_2^\times/\Q_2^{\times2}$ and $G'=G\langle-1\rangle$, so that $F_v=\Q_2(\sqrt G)$ and $F_v(i)=\Q_2(\sqrt{G'})$.
For a subgroup $H$ the ramification index of $\Q_2(\sqrt H)/\Q_2$ is $|H|/|H\cap\langle5\rangle|$. If $-1\in G$ the
extension $F_v(i)/F_v$ is trivial. Otherwise $|G'|=2|G|$, and $F_v(i)/F_v$ is unramified exactly when
$|G'\cap\langle5\rangle|=2|G\cap\langle5\rangle|$, that is, when $5\in G'\setminus G$, that is, when $-5\in G$. So
no prime above $2$ ramifies exactly when $G$ contains $-1$ or $-5$, that is, when the squarefree representative
$c\neq1$ of some element of the group generated by the $a_j$ modulo squares is $\equiv7$ or $\equiv3\pmod8$.
These $c$ are those of the quadratic subfields.
\end{proof}

For real quadratic fields Theorem~\ref{thm:main} is Fischer's theorem \cite[Theorem~8]{Fischer1990}. For $F$ of odd
degree, some prime $v$ above $2$ has odd degree over $\Q_2$; then $i\notin F_v$, and since $\Q_2(i)/\Q_2$ has
ramification index $2$ and $F_v/\Q_2$ an odd one, $F_v(i)/F_v$ is ramified, which is Moorhouse's Theorem~7.1
\cite{Moorhouse}.

\begin{remark}\label{rem:forms}
Nothing in the proof uses the particular form $x^2+y^2$. Let $c\in F^\times$ and $q=x^2+cy^2$. With
$L=F[t]/(t^2+c)$ and $\N(x+ty)=x^2+cy^2$, the same argument shows that the graph $q(x-x',y-y')=1$ on $F^2$ is
bipartite if and only if some prime of $F$ above $2$ ramifies in $F(\sqrt{-c})$. Every nondegenerate binary form
that represents $1$ is equivalent to such a form, and a form that does not represent $1$ gives a graph without
edges. Nondegeneracy is needed: for $q=x^2$ the components are the sets $(x_0+\Z)\times F$, and the graph is
bipartite over every $F$.
\end{remark}

\begin{example}\label{ex:quartic}
Theorem~\ref{thm:main} also decides fields that are not multiquadratic. Let $F$ be generated by a root of
$\alpha^4-2\alpha^3-2\alpha-3$ (Galois group $S_4$, so that $F$ has no quadratic subfield), of
$\alpha^4-6\alpha^2+4\alpha+2$ (whose only quadratic subfield is $\Q(\sqrt2)$) or of $\alpha^4-5\alpha^2-7$ (whose
only quadratic subfield is $\Q(\sqrt{53})$). In each case no prime above $2$ ramifies in $F(i)$, so $F^2$ has an
odd cycle, although the planes of $\Q(\sqrt2)$ and $\Q(\sqrt{53})$ are $2$-colourable. Exact integer relations
$\sum_k c_ku_k=0$ among unit vectors $u_k=z_k/\bar z_k$ with $\sum_k c_k$ odd, which are closed walks of odd length
$\sum_k|c_k|$, are found in these three fields by the script \texttt{odd\_walks.gp} of \cite{Repo}; for three
fields in which a prime above $2$ ramifies, the same search finds only relations with even sum. The script
\texttt{two\_colour\_criterion.gp} checks the criterion against Fischer's theorem for the $242$ squarefree $N$
with $2\le N\le400$, against Moorhouse's theorem for nine fields of degree $3$, $5$ and $7$, and against
Corollary~\ref{cor:multi} for $1\,522$ biquadratic fields.
\end{example}

\section{A local--global question}\label{sec:lg}

Call a number field \emph{real} if it has a real place. For a place $v$ of a real number field $F$, let
$\chiloc(v)$ be the least number of colours of a proper colouring of $F_v^2$ that is locally constant or, at an
archimedean place, Borel. At a real place this is the Borel chromatic number $\chiB(\R^2)$, which is at least $5$
by Falconer's theorem \cite{Falconer} and at most $7$ by the hexagonal colouring. At a complex place, and at a
finite place where $-1$ is a square in $F_v$, the form $x^2+y^2$ is isotropic and even the ordinary chromatic
number of $F_v^2$ is infinite, by a theorem of Davies \cite{Davies} (see \cite[Proposition~3]{PP}); there
$\chiloc(v)=\infty$. At a finite place where $x^2+y^2$ is anisotropic, the unit vectors are integral, and
$\chiloc(v)$ is the least chromatic number of the finite graphs $\mathrm{Cay}(\Oo_w/\m_w^r,\ T_r)$, where $w$ is
the place of $F(i)$ above $v$ and $T_r$ is the image of the norm-one units: unit steps preserve the cosets of
$\Oo_w$, which can be coloured independently, and a locally constant colouring of the compact set $\Oo_w$ is
constant on the cosets of some $\m_w^r$. Restriction gives
\[
\chi(F^2)\le\min_v\chiloc(v),
\]
and every upper bound below $7$ that we know for a plane over a number field is of this kind: reductions at one
finite place \cite{Fischer1990,Madore,MM,Moorhouse}.

\begin{question}\label{q:lg}
Is $\chi(F^2)=\min_v\chiloc(v)$ for every real number field $F$?
\end{question}

Theorem~\ref{thm:main} answers Question~\ref{q:lg} whenever one side is $2$: if $\chi(F^2)=2$,
Lemma~\ref{lem:ring} and Proposition~\ref{prop:onlyif} give a prime $v$ above $2$ that ramifies in $F(i)$, and
Proposition~\ref{prop:if} gives $\chiloc(v)=2$.

Two facts bound the local side from below. First, if $v$ lies above a prime $p$, then $F_v\supseteq\Q_p$, and a locally
constant colouring of $F_v^2$ restricts to a locally constant, hence measurable, colouring of $\Q_p^2$; so
$\chiloc(v)\ge\chiM(\Q_p^2)$, the measurable chromatic number. For $p\equiv3\pmod4$, the ratio bound of Delsarte
and Hoffman on the compact group $\Z_p[i]$ gives $\chiM(\Q_p^2)\ge4$ for $p\ge11$, $\chiM(\Q_p^2)\ge5$ for
$p\ge23$, and $\chiM(\Q_p^2)\ge7$ for $p=71$ and for $p\ge103$ \cite[Theorem~1 and Table~1]{PP}. Second, if a
field $K$ embeds in $F_v$, then $\chiloc(v)\ge\chi(K^2)$.

\textbf{Three colours.} For squarefree $d\ge2$, the field $\Q(\sqrt d)$ has a place with $\chiloc(v)\le3$ exactly
when $d\not\equiv11\pmod{12}$. If $d\not\equiv3\pmod4$, Proposition~\ref{prop:if} gives two colours at $2$. If
$d\not\equiv2\pmod3$, a place $v$ above $3$ has residue field $\F_3$, the place above it has residue field $\F_9$,
and the norm-one units reduce to $\{\pm1,\pm i\}$; so $a+bi\mapsto a+b\bmod3$ colours the first level with three
colours (cf.\ \cite[Proposition~3.2]{Madore}, \cite[Theorem~9]{Fischer1990}). Conversely, let
$d\equiv11\pmod{12}$. The real places have $\chiloc\ge5$. The prime $3$ is inert, and the places above $3$, the
places above primes $p\equiv1\pmod4$, and the place above $2$ when $d\equiv7\pmod8$ contain $i$. When
$d\equiv3\pmod8$, the place above $2$ has $F_v=\Q_2(\sqrt3)$, which contains $\sqrt{11}$ because $33$ is a square
in $\Q_2$; and a place above $7$ has $F_v\supseteq\Q_7$, which contains $\sqrt{11}$ because $11\equiv2^2\pmod7$.
In both cases $\chiloc(v)\ge\chi(\Q(\sqrt{11})^2)=4$ \cite{QP}. Above a prime $p\equiv3\pmod4$ with $p\ge11$,
$\chiloc(v)\ge\chiM(\Q_p^2)\ge4$. So Question~\ref{q:lg} predicts
$\chi(\Q(\sqrt d)^2)\ge4$ for every squarefree $d\equiv11\pmod{12}$. This has been checked, by unit-distance
graphs with certified proofs of non-$3$-colourability, for $27$ values of $d$ \cite{QP}; no counterexample is
known.

Colourings that are locally constant at several anisotropic finite places colour a categorical product of finite
level graphs, since $T$ is dense in the product of the local groups $T(F_v)$; they can do better than every
single place only through a counterexample to Hedetniemi's conjecture, as Moorhouse observed
\cite[\S9]{Moorhouse}. Such counterexamples exist with many colours \cite{Shitov}, so for large numbers of colours
Question~\ref{q:lg} is a genuine restriction to single places; with three colours they do not exist \cite{EZS}.

\textbf{Four colours.} Call a number field $F$ \emph{admissible} if $\sqrt3\notin F$ and $i\in F_v$ for every
place $v$ of $F$ above $2$, $3$, $7$, $11$ and $19$. A real admissible field has $\chiloc(v)\ge5$ at every
place $v$: $\chiloc\ge5$ at the real places; $\chiloc=\infty$ at the complex places, at the
places above $2$, $3$, $7$, $11$ and $19$ and at those above primes $p\equiv1\pmod4$; and at a place above a
prime $p\equiv3\pmod4$ with $p\ge23$, $\chiloc(v)\ge\chiM(\Q_p^2)\ge5$. The margins are small: the ratio bound
on the density of a colour class is $0.24908$ at $p=23$ and $0.2499973$ at $p=31$, against $1/4$
\cite[Table~1]{PP}. An admissible field also has no unit triangle, since the third vertex of a unit equilateral
triangle on $0$ and $(a,b)$ is $(\frac a2-\frac{\sqrt3}2b,\frac b2+\frac{\sqrt3}2a)$. The smallest admissible real
quadratic field is $\Q(\sqrt{167})$; the next are $d=887$, $1055$, $1319$ and $1823$. Fischer had singled out the
same field: $\Q(\sqrt{167})^2$ has no additive $k$-colouring for $k<11$ \cite[Example~4]{Fischer1990}. So a positive answer to
Question~\ref{q:lg} for $\Q(\sqrt{167})$ would give $\chi(\Q(\sqrt{167})^2)\ge5$, and by the theorem of de Bruijn
and Erd\H{o}s \cite{dBE} a finite triangle-free $5$-chromatic unit-distance graph in the plane. According to the
abstract of \cite[Chapter~55]{Soifer2}, Soifer asked for such graphs in a problem paper of 2019; according to
\cite{MW}, de Grey has found one in $\R^3$. We have not found a planar one in the literature we could consult. The searches of \cite{Repo} have not
found one over $\Q(\sqrt{167})$: its unit vectors have large denominators, and the graphs grown there stay sparse.

\textbf{The real place.} If Question~\ref{q:lg} had a positive answer for every real quadratic field, then
$\chi(\R^2)=\chiB(\R^2)$. Take a prime $d\equiv7\pmod8$ that is a quadratic non-residue modulo every prime
$p\equiv3\pmod4$ below $100$ other than $71$; the smallest is $d=186\,023$. In $\Q(\sqrt d)$ the places above
$2$, above the primes $p\equiv1\pmod4$ and above those primes $p\equiv3\pmod4$ contain $i$, and every other finite
place lies above $71$ or above a prime $p\ge103$, so $\chiloc(v)\ge\chiM(\Q_p^2)\ge7$. The question would then give
$\chi(\Q(\sqrt d)^2)=\chiB(\R^2)$, while $\chi(\Q(\sqrt d)^2)\le\chi(\R^2)\le\chiB(\R^2)$. Whether
$\chi(\R^2)=\chiB(\R^2)$ is open, so this part of the question is the most speculative.

\begin{thebibliography}{99}
\bibitem{AM} M. F. Atiyah and I. G. Macdonald, \emph{Introduction to commutative algebra}, Addison-Wesley,
Reading, Mass., 1969.
\bibitem{BP} M. Benda and M. Perles, \emph{Colorings of metric spaces}, Geombinatorics \textbf{9} (2000),
113--126 (not seen by us).
\bibitem{Davies} J. Davies, \emph{Chromatic number of spacetime}, arXiv:2308.16885 (2023, revised 2024).
\bibitem{dBE} N. G. de Bruijn and P. Erd\H{o}s, \emph{A colour problem for infinite graphs and a problem in the
theory of relations}, Nederl. Akad. Wetensch. Proc. Ser. A \textbf{54} = Indag. Math. \textbf{13} (1951),
369--373.
\bibitem{deGrey} A. D. N. J. de Grey, \emph{The chromatic number of the plane is at least 5}, Geombinatorics
\textbf{28} (2018), 18--31; arXiv:1804.02385.
\bibitem{EZS} M. El-Zahar and N. Sauer, \emph{The chromatic number of the product of two $4$-chromatic graphs is
$4$}, Combinatorica \textbf{5} (1985), 121--126.
\bibitem{EI} G. Exoo and D. Ismailescu, \emph{The chromatic number of the plane is at least $5$: a new proof},
Discrete Comput. Geom. \textbf{64} (2020), 216--226.
\bibitem{Falconer} K. J. Falconer, \emph{The realization of distances in measurable subsets covering $\R^n$},
J. Combin. Theory Ser. A \textbf{31} (1981), 184--189.
\bibitem{Fischer1990} K. G. Fischer, \emph{Additive $K$-colorable extensions of the rational plane}, Discrete
Math. \textbf{82} (1990), 181--195.
\bibitem{Repo} S. Gal\'an, \emph{The Hadwiger--Nelson problem over number fields}, code, data and formal
proofs, 2026. Archived at Zenodo, \href{https://doi.org/10.5281/zenodo.22976635}{doi:10.5281/zenodo.22976635}
(all versions); source at \url{https://github.com/decalion89/chromatic-number-of-the-plane}; the scripts cited here
are in \texttt{data/quadratic\_planes/scripts/}, and the account is \texttt{notes/local\_global.md}.
\bibitem{PP} S. Gal\'an, \emph{Colouring the $p$-adic plane}, draft, 2026, in \cite{Repo}
(\texttt{papers/padic-planes/}).
\bibitem{QP} S. Gal\'an, \emph{Real quadratic planes that need four colours}, draft, 2026, in \cite{Repo}
(\texttt{papers/quadratic-planes/}).
\bibitem{Johnson1987} P. D. Johnson Jr., \emph{Two-colorings of real quadratic extensions of $\Q^2$ that forbid
many distances}, Congr. Numer. \textbf{60} (1987), 51--58 (not seen by us).
\bibitem{Johnson2006} P. D. Johnson Jr., \emph{Coloring the rational points to forbid the distance one---a
tentative history and compendium}, Geombinatorics \textbf{16} (2006), 209--218 (not seen by us).
\bibitem{Madore} D. A. Madore, \emph{The Hadwiger--Nelson problem over certain fields}, arXiv:1509.07023 (2015).
\bibitem{MM} MildlyMeticulous (GitHub account), \emph{A $2$-adic obstruction to $5$-chromatic unit-distance
graphs, and the exact chromatic number of the plane over $\Q(\sqrt3,\sqrt{11})$}, code, logs and draft, July 2026,
\url{https://github.com/MildlyMeticulous/hn-2adic-obstruction} (file \texttt{RESULTS.md}, \S1).
\bibitem{Moorhouse} G. E. Moorhouse, \emph{On the chromatic numbers of planes}, preliminary draft, revised
3 March 2010, \url{https://www.ericmoorhouse.org/pub/chromatic.pdf}.
\bibitem{Payne} M. S. Payne, \emph{Unit distance graphs with ambiguous chromatic number}, Electron. J. Combin.
\textbf{16} (2009), Note 31; arXiv:0707.1177.
\bibitem{Shitov} Y. Shitov, \emph{Counterexamples to Hedetniemi's conjecture}, Ann. of Math. (2) \textbf{190}
(2019), 663--667.
\bibitem{Soifer} A. Soifer, \emph{The mathematical coloring book}, Springer, New York, 2009.
\bibitem{Soifer2} A. Soifer, \emph{The new mathematical coloring book}, 2nd ed., Springer, New York, 2024
(we have read only the abstract of Chapter~55, ``Triangle-free 5-chromatic unit distance graphs'').
\bibitem{MW} E. W. Weisstein, \emph{Gr\"otzsch graph}, MathWorld, \url{https://mathworld.wolfram.com/GroetzschGraph.html},
citing A. D. N. J. de Grey, \emph{A 5-chromatic, triangle-free unit-distance graph in $\R^3$ with 61 vertices},
Geombinatorics \textbf{35} (2026).
\bibitem{Woodall} D. R. Woodall, \emph{Distances realized by sets covering the plane}, J. Combin. Theory Ser. A
\textbf{14} (1973), 187--200.
\end{thebibliography}

\end{document}
